\documentclass[11pt]{amsart}
\usepackage{geometry}
\usepackage{graphicx}
\usepackage{amssymb}
\usepackage{epstopdf}

\geometry{a4paper}

\newtheorem*{theorem*}{Theorem}
\newtheorem*{lemma*}{Lemma}
\newtheorem*{proposition*}{Proposition}

\title{\(\sqrt{2}\) is not a rational number}
\author{Minqi Pan}
\date{\today}

\begin{document}
\maketitle

\begin{lemma*}
$\forall n\in \mathbb{N}$, $n^2$ is even implies that $n$ is even.
\end{lemma*}

\begin{proof}
Suppose that $n$ is odd, then
\[n=2k+1\]
for some $k\in \mathbb{N}$. Thus,
\[n^2=4k^2+4k+1=2(2k^2+2k)+1\]
which means $n^2$ is odd. The lemma thus holds by contraposition.
\end{proof}

\begin{proposition*}
\(\sqrt{2}\) is not a rational number.
\end{proposition*}

\begin{proof}
Suppose that \(\sqrt{2}\) is a rational number, then
\[\sqrt{2}=\frac{m}{n}\]
for some $m\in \mathbb{N}$, $n\in \mathbb{N}$ and $(m, n)$ are co-prime integers. Also,
\[2=\frac{m^2}{n^2}\]
which means that
\begin{equation} \label{1}
m^2=2n^2
\end{equation}
thus $m^2$ is even. And by Lemma, $m$ is even. Thus there exists $k\in \mathbb{N}$ such that
\begin{equation} \label{2}
m=2k
\end{equation}
Now substitute equation~\ref{2} into equation~\ref{1},
\[
4k^2=2n^2
\]
That is,
\[
2k^2=n^2
\]
Thus $n^2$ is even. By Lemma, $n$ is even. Thus $n$ and $m$ are both even, contradicting the proposition that $(m, n)$ are co-prime integers.
\end{proof}

\end{document}
